% Supplementary material for
%   "Declared Ambiguity: Fail-Closed Execution for Non-Idempotent Legacy APIs
%    Without Idempotency Keys"
%
% WHY THIS IS A SEPARATE DOCUMENT AND NOT AN APPENDIX TO main.tex.
%
% The IEEE Computer Society author guidance, which governs TSE, states that
% "appendices in journal articles are treated as supplemental material and must
% be submitted separately rather than in the main manuscript file", and that
% "Supplemental material is not included in these word counts". Both sentences
% point the same way: an appendix bound into main.tex would still be part of the
% manuscript file, would still count against the 12-page regular-paper limit,
% and would have to be pulled back out at submission. A separate PDF is what the
% submission process asks for, so that is what this is. Convenience did not
% decide it -- if anything an appendix would have been easier to build.
%
% CROSS-REFERENCES. Being a separate document, this one cannot \cref labels
% defined in main.tex. References to the paper are therefore by name, spelled
% out, so a reader holding only this file can still find the passage. That is a
% deliberate cost of the separation and not an oversight.
%
% It carries the same anonymity toggle as main.tex, and build_paper.sh applies
% the same pdfTeX primitives to its anonymous build, because a second PDF that
% leaks the author is the same failure as the first one leaking.

\documentclass[10pt,journal,compsoc]{IEEEtran}

\usepackage[T1]{fontenc}
\usepackage{amsmath}
\usepackage{booktabs}
\usepackage{tabularx}
\usepackage{graphicx}
\usepackage[hidelinks]{hyperref}
\usepackage[capitalise,noabbrev]{cleveref}

% The same generated macros the paper uses. Numbers here are generated by
% scripts/paper_tables.py exactly as the paper's are; none is typed by hand,
% and check_paper_numbers.py scans this file so a macro used only here is not
% reported orphaned.
% GENERATED by scripts/paper_tables.py -- do not edit.
% Sources: analysis/per-cell-metrics.csv, analysis/latency-and-throughput.csv,
%          analysis/redis-kill-ablation.csv, analysis/comparisons-vs-aep-full.csv,
%          experiments/results/g2-flakey-write-loss*.json
% Each macro carries the file, filter and arithmetic behind it.

% per-cell-metrics.csv | system=AEP_FULL regime=(session-3) response_class=AUTHORITATIVE_READBACK
% metric=known_ambiguity_rate | sum(successes)/sum(total) = 0/180
\newcommand{\AepAmbAuth}{0.0000}

% per-cell-metrics.csv | system=AEP_FULL regime=(session-3) response_class=AUTHORITATIVE_READBACK
% metric=undetected_duplicate_rate | sum(successes)/sum(total) = 0/180
\newcommand{\AepDupAuth}{0.0000}

% per-cell-metrics.csv | system=AEP_FULL regime=(session-3) response_class=AUTHORITATIVE_READBACK
% metric=lost_effect_rate | sum(successes)/sum(total) = 0/180
\newcommand{\AepLostAuth}{0.0000}

% per-cell-metrics.csv | system=AEP_FULL regime=(session-3) response_class=POSITIVE_ONLY_READBACK
% metric=known_ambiguity_rate | sum(successes)/sum(total) = 63/180
\newcommand{\AepAmbPosOnly}{0.3500}

% per-cell-metrics.csv | system=AEP_FULL regime=(session-3) response_class=POSITIVE_ONLY_READBACK
% metric=undetected_duplicate_rate | sum(successes)/sum(total) = 0/180
\newcommand{\AepDupPosOnly}{0.0000}

% per-cell-metrics.csv | system=AEP_FULL regime=(session-3) response_class=POSITIVE_ONLY_READBACK
% metric=lost_effect_rate | sum(successes)/sum(total) = 0/180
\newcommand{\AepLostPosOnly}{0.0000}

% per-cell-metrics.csv | system=AEP_FULL regime=(session-3) response_class=NO_READBACK
% metric=known_ambiguity_rate | sum(successes)/sum(total) = 130/180
\newcommand{\AepAmbNoReadback}{0.7222}

% per-cell-metrics.csv | system=AEP_FULL regime=(session-3) response_class=NO_READBACK
% metric=undetected_duplicate_rate | sum(successes)/sum(total) = 0/180
\newcommand{\AepDupNoReadback}{0.0000}

% per-cell-metrics.csv | system=AEP_FULL regime=(session-3) response_class=NO_READBACK
% metric=lost_effect_rate | sum(successes)/sum(total) = 0/180
\newcommand{\AepLostNoReadback}{0.0000}

% per-cell-metrics.csv | executions behind every AEP-full rate on this capability class (regime=(session-3))
\newcommand{\AepExecAuth}{180}

% per-cell-metrics.csv | executions behind every AEP-full rate on this capability class (regime=(session-3))
\newcommand{\AepExecPosOnly}{180}

% per-cell-metrics.csv | executions behind every AEP-full rate on this capability class (regime=(session-3))
\newcommand{\AepExecNoReadback}{180}

% per-cell-metrics.csv | min undetected-duplicate rate over {B0,B1,B2} x {all collected response classes}
% regime=(session-3) | 9 cells
\newcommand{\BaselineDupLow}{0.77}

% per-cell-metrics.csv | max undetected-duplicate rate over {B0,B1,B2} x {all collected response classes}
% regime=(session-3) | 9 cells
\newcommand{\BaselineDupHigh}{0.83}

% \BaselineDupLow as a percentage, for comparison against \AblationZeroUpper
\newcommand{\BaselineDupLowPct}{77}

% \BaselineDupHigh as a percentage, for comparison against \AblationZeroUpper
\newcommand{\BaselineDupHighPct}{83}

% comparisons-vs-aep-full.csv | metric=undetected_duplicate_rate, largest (weakest) Fisher p over {B0,B1,B2} vs AEP_FULL
% 3 comparisons; the other 2 are smaller
\newcommand{\BaselineDupMaxP}{2.1\times10^{-183}}

% latency-and-throughput.csv | system=AEP_FULL
% step_latency_ms_median over crash-free runs only (overhead_runs_crash_free=3)
% E5 gate: runs without a declared suspend-disabled host contribute NO timing
\newcommand{\AepStepMedian}{4\,004.9}

% latency-and-throughput.csv | system=B0_NAIVE_RETRY
% step_latency_ms_median over crash-free runs only (overhead_runs_crash_free=3)
% E5 gate: runs without a declared suspend-disabled host contribute NO timing
\newcommand{\BzeroStepMedian}{2\,010.2}

% latency-and-throughput.csv | system=B3_INTENT_NO_BARRIER
% step_latency_ms_median over crash-free runs only (overhead_runs_crash_free=3)
% E5 gate: runs without a declared suspend-disabled host contribute NO timing
\newcommand{\BthreeStepMedian}{2\,038.2}

% latency-and-throughput.csv | B3 median - B0 median
% = 2038.2 - 2010.2; the whole write-ahead protocol except the barrier
\newcommand{\ProtocolMinusBarrier}{28.0}

% = (2038.2 - 2010.2) / 2010.2 x 100
\newcommand{\ProtocolMinusBarrierPct}{1.4}

% latency-and-throughput.csv | AEP-full median - B3 median, both under appendfsync=everysec
% = 4004.9 - 2038.2; the two WAITAOF round trips together
\newcommand{\BarrierCost}{1\,966.7}

% half of \BarrierCost -- the protocol runs exactly two barriers per step
\newcommand{\BarrierCostEach}{983.3}

% = \BarrierCostEach / \AepStepMedian x 100
% = 983.3 / 4004.9 x 100; the end-to-end cost of adding one more barrier to the step, not the increase in the barrier bill alone
\newcommand{\ThirdBarrierStepPct}{24.6}

% latency-and-throughput.csv | \BarrierCost / \ProtocolMinusBarrier, both under appendfsync=everysec
% = (4004.9 - 2038.2) / (2038.2 - 2010.2) = 1966.7 / 28.0 = 70.14, to two significant figures
% how much more the two fsync barriers cost than everything else the protocol does
\newcommand{\BarrierToProtocolRatio}{70}

% fsync-always/analysis/latency-and-throughput.csv | AEP-full median - B3 median, both under appendfsync=always
% = 2063.4 - 2048.4
\newcommand{\BarrierCostAlways}{15.0}

% fsync-always/analysis/latency-and-throughput.csv | system=B3_INTENT_NO_BARRIER, appendfsync=always
% against 2038.2 ms under everysec: the ablated protocol's own cost is nearly policy-independent, which is what makes the two barrier figures comparable
\newcommand{\BthreeAlwaysMedian}{2\,048.4}

% fsync-always/analysis/latency-and-throughput.csv | system=AEP_FULL, appendfsync=always
% step_latency_ms_median; quoted in the threats section against the p95 below, to show the tail the three-run interval is drawn from
\newcommand{\AepAlwaysMedian}{2\,063.4}

% fsync-always/analysis/latency-and-throughput.csv | system=AEP_FULL, appendfsync=always
% step_latency_ms_p95
\newcommand{\AepAlwaysPninetyfive}{5\,067.8}

% fsync-always/analysis/latency-and-throughput.csv | system=AEP_FULL, appendfsync=always
% executions_per_second
\newcommand{\AepThroughputAlways}{0.42}

% latency-and-throughput.csv | system=AEP_FULL, appendfsync=everysec
% executions_per_second; the everysec matrix collected more runs per cell, so wall time per execution is not comparable as a benchmark -- it is quoted only against the always arm
\newcommand{\AepThroughputEverysec}{0.33}

% per-cell-metrics.csv | system=B4_DURABLE_WORKFLOW regime=(session-3) response_class=AUTHORITATIVE_READBACK
% metric=undetected_duplicate_rate | sum(successes)/sum(total) = 95/180
\newcommand{\BfourDupAuth}{0.5278}

% per-cell-metrics.csv | executions behind \BfourDupAuth
\newcommand{\BfourExecAuth}{180}

% per-cell-metrics.csv | system=B4_DURABLE_WORKFLOW regime=(session-3) response_class=POSITIVE_ONLY_READBACK
% metric=undetected_duplicate_rate | sum(successes)/sum(total) = 106/180
\newcommand{\BfourDupPosOnly}{0.5889}

% per-cell-metrics.csv | executions behind \BfourDupPosOnly
\newcommand{\BfourExecPosOnly}{180}

% per-cell-metrics.csv | system=B4_DURABLE_WORKFLOW regime=(session-3) response_class=NO_READBACK
% metric=undetected_duplicate_rate | sum(successes)/sum(total) = 101/180
\newcommand{\BfourDupNoReadback}{0.5611}

% per-cell-metrics.csv | executions behind \BfourDupNoReadback
\newcommand{\BfourExecNoReadback}{180}

% per-cell-metrics.csv | system=B4B_DURABLE_WORKFLOW_AT_MOST_ONCE regime=(session-3) response_class=AUTHORITATIVE_READBACK
% metric=lost_effect_rate | sum(successes)/sum(total) = 98/180
\newcommand{\BfourbLostAuth}{0.5444}

% per-cell-metrics.csv | executions behind \BfourbLostAuth
\newcommand{\BfourbExecAuth}{180}

% per-cell-metrics.csv | system=B4B_DURABLE_WORKFLOW_AT_MOST_ONCE regime=(session-3) response_class=POSITIVE_ONLY_READBACK
% metric=lost_effect_rate | sum(successes)/sum(total) = 91/180
\newcommand{\BfourbLostPosOnly}{0.5056}

% per-cell-metrics.csv | executions behind \BfourbLostPosOnly
\newcommand{\BfourbExecPosOnly}{180}

% per-cell-metrics.csv | system=B4B_DURABLE_WORKFLOW_AT_MOST_ONCE regime=(session-3) response_class=NO_READBACK
% metric=lost_effect_rate | sum(successes)/sum(total) = 93/180
\newcommand{\BfourbLostNoReadback}{0.5167}

% per-cell-metrics.csv | executions behind \BfourbLostNoReadback
\newcommand{\BfourbExecNoReadback}{180}

% per-cell-metrics.csv | max known_ambiguity_rate over {B4,B4b} x {all capability classes} x {all crash points}
% regime=(session-3) | 36 cells
% the durable-execution engine never declares ambiguity; that is the corner of the trilemma it cannot reach
\newcommand{\BfourFamilyAmbMax}{0.0000}

% per-cell-metrics.csv | cells behind \BfourFamilyAmbMax
\newcommand{\BfourFamilyAmbCells}{36}

% per-cell-metrics.csv | system=B3_INTENT_NO_BARRIER regime=(session-3) response_class=AUTHORITATIVE_READBACK
% metric=known_ambiguity_rate | sum(successes)/sum(total) = 0/180
\newcommand{\BthreeAmbAuth}{0.0000}

% per-cell-metrics.csv | system=B3_INTENT_NO_BARRIER regime=(session-3) response_class=POSITIVE_ONLY_READBACK
% metric=known_ambiguity_rate | sum(successes)/sum(total) = 66/180
\newcommand{\BthreeAmbPosOnly}{0.3667}

% per-cell-metrics.csv | system=B3_INTENT_NO_BARRIER regime=(session-3) response_class=NO_READBACK
% metric=known_ambiguity_rate | sum(successes)/sum(total) = 129/180
\newcommand{\BthreeAmbNoReadback}{0.7167}

% comparisons-vs-aep-full.csv | metric=undetected_duplicate_rate system=B3_INTENT_NO_BARRIER reference=AEP_FULL
% B3 0/600 vs AEP-full 0/600
% the ablation is indistinguishable from the full protocol here; that is the finding, not a caveat
\newcommand{\BthreeVsAepDupP}{1.00}

% comparisons-vs-aep-full.csv | metric=lost_effect_rate system=B3_INTENT_NO_BARRIER reference=AEP_FULL
% B3 0/600 vs AEP-full 0/600
% the ablation is indistinguishable from the full protocol here; that is the finding, not a caveat
\newcommand{\BthreeVsAepLostP}{1.00}

% comparisons-vs-aep-full.csv | metric=known_ambiguity_rate system=B3_INTENT_NO_BARRIER reference=AEP_FULL
% B3 225/600 vs AEP-full 223/600
% the ablation is indistinguishable from the full protocol here; that is the finding, not a caveat
\newcommand{\BthreeVsAepAmbP}{0.95}

% comparisons-vs-aep-full.csv | executions per arm, identical across all three metrics
\newcommand{\BthreeVsAepN}{600}

% Wilson 95% upper limit on 0/600, as a percentage
% the largest undetected-duplicate or lost-effect rate either system could have and still record zero; the difference between them is bounded by it
\newcommand{\AblationZeroUpper}{0.64}

% comparisons-vs-aep-full.csv | metric=undetected_duplicate_rate, the larger of the two arms' numerators
% B3 0/600 vs AEP-full 0/600
\newcommand{\BthreeVsAepDupCount}{0}

% comparisons-vs-aep-full.csv | metric=lost_effect_rate, the larger of the two arms' numerators
% B3 0/600 vs AEP-full 0/600
\newcommand{\BthreeVsAepLostCount}{0}

% comparisons-vs-aep-full.csv | metric=known_ambiguity_rate, |B3 successes - AEP-full successes|
% |225 - 223| over 600 executions per arm
\newcommand{\BthreeVsAepAmbDelta}{2}

% redis-kill-ablation.csv | regime=redis-kill-preack system=AEP_FULL response_class=NO_READBACK
% executions_with_an_applied_effect
\newcommand{\AepKillApplied}{10}

% redis-kill-ablation.csv | regime=redis-kill-preack system=AEP_FULL response_class=NO_READBACK
% unwanted-applied-effect rate = executions_with_an_applied_effect/executions = 10/30
% an effect put on the wire while the durability of its own intent record could no longer be confirmed
\newcommand{\AepUnwantedRate}{0.3333}

% redis-kill-ablation.csv | runs for AEP_FULL
\newcommand{\AepKillRuns}{30}

% redis-kill-ablation.csv | canary_survived/(survived+lost) for AEP_FULL
% the un-acknowledged write made immediately before the hard kill
\newcommand{\AepKillCanary}{30/30}

% redis-kill-ablation.csv | regime=redis-kill-preack system=B3_INTENT_NO_BARRIER response_class=NO_READBACK
% executions_with_an_applied_effect
\newcommand{\BthreeKillApplied}{28}

% redis-kill-ablation.csv | regime=redis-kill-preack system=B3_INTENT_NO_BARRIER response_class=NO_READBACK
% unwanted-applied-effect rate = executions_with_an_applied_effect/executions = 28/30
% an effect put on the wire while the durability of its own intent record could no longer be confirmed
\newcommand{\BthreeUnwantedRate}{0.9333}

% redis-kill-ablation.csv | runs for B3_INTENT_NO_BARRIER
\newcommand{\BthreeKillRuns}{30}

% redis-kill-ablation.csv | canary_survived/(survived+lost) for B3_INTENT_NO_BARRIER
% the un-acknowledged write made immediately before the hard kill
\newcommand{\BthreeKillCanary}{30/30}

% redis-kill-ablation.csv | B3 applied - AEP-full applied
% = 28 - 10; real non-idempotent effects that the barrier withheld under an identical injected fault
\newcommand{\UnwantedPrevented}{18}

% redis-kill-ablation.csv | Fisher exact two-tailed on [[10, 20], [28, 2]]
% AEP-full vs B3 on the unwanted-applied-effect rate
\newcommand{\UnwantedP}{1.9\times10^{-6}}

% redis-kill-ablation.csv | canary_survived + canary_lost, summed over AEP_FULL and B3_INTENT_NO_BARRIER
% un-acknowledged writes made immediately before a hard Redis kill, inside the ablation cells
\newcommand{\KillCanaryN}{60}

% experiments/results/g2-flakey-write-loss*.json | 3 replication(s) of experiments/flakey_write_loss.py
% independent runs of the probe, each rebuilding the device stack and the filesystem from scratch
\newcommand{\FlakeyReplications}{3}

% experiments/results/g2-flakey-write-loss*.json | 3 replication(s) of experiments/flakey_write_loss.py
% countable trials: the acknowledged write survived, so the trial says something about the unacknowledged one
\newcommand{\FlakeyN}{90}

% experiments/results/g2-flakey-write-loss*.json | 3 replication(s) of experiments/flakey_write_loss.py
% WAITAOF-acknowledged writes still present after the device stopped accepting writes
\newcommand{\FlakeyAckSurvived}{90/90}

% experiments/results/g2-flakey-write-loss*.json | 3 replication(s) of experiments/flakey_write_loss.py
% un-acknowledged writes destroyed by the same event
\newcommand{\FlakeyUnackLost}{90/90}

% experiments/results/g2-flakey-write-loss*.json | 3 replication(s) of experiments/flakey_write_loss.py
% narrowest write-to-write-loss exposure window, ms
\newcommand{\FlakeyWindowMin}{12.4}

% experiments/results/g2-flakey-write-loss*.json | 3 replication(s) of experiments/flakey_write_loss.py
% widest write-to-write-loss exposure window, ms; the appendfsync everysec period it sits inside is 1000 ms
\newcommand{\FlakeyWindowMax}{61.8}

% experiments/results/g2-flakey-write-loss*.json | 3 replication(s) of experiments/flakey_write_loss.py
% Fisher exact two-tailed, acknowledged vs un-acknowledged loss over the same 90 trials
\newcommand{\FlakeyBarrierP}{2.2\times10^{-53}}

% reports/raw/e1-durability-window.txt (0/10 lost under docker kill -s KILL) vs experiments/results/g2-flakey-write-loss*.json | 3 replication(s) of experiments/flakey_write_loss.py
% Fisher exact two-tailed across the two fault classes
\newcommand{\FlakeyVsProcessKillP}{5.8\times10^{-14}}

% analysis/per-execution.csv | cluster bootstrap over runs, 10000 resamples, seed 20260806, appendfsync=everysec
% 2.5th percentile of (median AEP-full - median B3); point estimate 1968.1 ms from 3 and 3 runs
\newcommand{\BarrierCostLow}{477.9}

% analysis/per-execution.csv | 97.5th percentile of the same bootstrap, appendfsync=everysec
\newcommand{\BarrierCostHigh}{1\,978.8}

% analysis/per-execution.csv | cluster bootstrap over runs, 10000 resamples, seed 20260806, appendfsync=always
% 2.5th percentile of (median AEP-full - median B3); point estimate 16.5 ms from 3 and 3 runs
\newcommand{\BarrierCostAlwaysLow}{-1\,474.6}

% analysis/per-execution.csv | 97.5th percentile of the same bootstrap, appendfsync=always
\newcommand{\BarrierCostAlwaysHigh}{22.7}

% analysis/coverage.json | runs
\newcommand{\RunsCollected}{432}

% analysis/coverage.json | executions
\newcommand{\ExecutionsCollected}{3\,780}

% analysis/coverage.json | cells
\newcommand{\CellsCollected}{126}

% analysis/coverage.json | bootstrap_resamples
\newcommand{\BootstrapResamples}{10\,000}

% reports/raw/e1-durability-window.txt | un-acknowledged writes lost to `docker kill -s KILL` under appendfsync=everysec
% everysec defers the fsync, not the write; a still-running kernel flushes the page cache
\newcommand{\ProcessKillUnackLost}{0/10}

% reports/raw/e1-durability-window.txt | usable trials
\newcommand{\ProcessKillTrials}{10}

% reports/raw/e1-durability-window.txt | min write->death over 10 trials, ms
\newcommand{\ProcessKillWindowMin}{419}

% reports/raw/e1-durability-window.txt | max write->death over 10 trials, ms
% every kill landed inside the 1000 ms fsync period, which is what makes the zero above a measurement rather than a miss
\newcommand{\ProcessKillWindowMax}{992}

% lines of Python under aep_core/, excluding __pycache__
% 14 files; regenerated on every run of this script
\newcommand{\CoreLoc}{6\,849}

% lines of Python under experiments/, excluding __pycache__
% 85 files; regenerated on every run of this script
\newcommand{\HarnessLoc}{21\,458}



% ---------------------------------------------------------------------------
% Anonymity toggle. Identical in shape to main.tex's: the flag is defined on
% the pdflatex command line by build_paper.sh --anonymous, so neither branch can
% be reached by editing this file alone, and an identity added to the wrong
% branch cannot reach supplementary-anon.pdf.
% ---------------------------------------------------------------------------
\newif\ifanonymous
\anonymousfalse
\ifdefined\ANONYMOUS \anonymoustrue \fi

\begin{document}

\title{Supplementary Material for\\
\emph{Declared Ambiguity: Fail-Closed Execution for Non-Idempotent Legacy
APIs Without Idempotency Keys}}

\ifanonymous
  % Mirrors main.tex's anonymous branch, including the DocInfo wipe. A second
  % PDF that leaks the author is the same failure as the first one leaking.
  \author{Anonymous Author(s)}
  \hypersetup{pdfauthor={},pdftitle={},pdfsubject={},pdfkeywords={},
              pdfcreator={},pdfproducer={}}
\else
  % Deliberately empty until the main paper's byline is settled, so the two
  % documents cannot disagree. What check_paper_numbers.py asserts is the leak
  % that actually matters: that the anonymous supplementary carries the
  % anonymous byline and does NOT contain the main paper's public byline.
  \author{}
\fi

\maketitle

\section{Purpose and status}

This document holds material referenced from the main paper but not required to
follow its argument: the enumerated coverage gaps of the fault-injection matrix,
and the measurement incidents found during collection. It is submitted as
supplemental material rather than as an appendix, for the reason recorded at the
top of this file.

Every number here is generated by the same script that generates the paper's,
from the same frozen results, and is checked by the same gate. Section names in
the main paper are given in words rather than as cross-references, because the
two documents are compiled separately.

\section{Coverage of the fault-injection matrix}
\label{supp:coverage}

The main paper's external-validity subsection states the headline coverage and
points here. The full plan is 1\,068 runs; we collected \RunsCollected{},
comprising \ExecutionsCollected{} executions over \CellsCollected{} cells,
chosen to answer the research questions rather than to fill the grid. The run
count for every cell is in \texttt{experiments/results/matrix/MANIFEST.md} and
in \texttt{analysis/coverage.json} in the artifact.

What that leaves uncovered, precisely --- six named gaps, of which three have
since been closed and are recorded here with what closing them changed:

\begin{itemize}
\item \textbf{B4's and B4b's \textsc{auth} cells were incomplete and are now
      collected --- and completing them corrected the value.} In an earlier
      draft B4's \textsc{auth} cell held twenty executions and B4b had none, so
      that cell printed as a dash. Both are now $n=\BfourExecAuth{}$ and
      $n=\BfourbExecAuth{}$ over all six crash points, the same shape as every
      other cell in the paper's outcomes table, and that table is free of
      dashes. We record what the completion changed, because it is a caution
      about partial cells generally and not only about this one: the earlier
      twenty executions came from a \emph{single} crash point, so the figure
      that draft quoted was not a wide interval around the value the full cell
      reports --- it was one crash point standing in for six, and the six differ
      enough that none of them represents the pooled rate. The full cell puts
      B4's \textsc{auth} duplicate rate at \BfourDupAuth{}, in line with its
      other two classes and far from the figure the partial cell suggested. The
      direction of the argument is unchanged, which is why we could describe the
      gap as non-essential before collecting it, but the number we quoted for it
      was wrong, and a partial cell being unrepresentative rather than merely
      imprecise is the more general lesson. Every other cell in the outcomes
      table spans all the crash points that apply to its system: six for the
      systems that run \texttt{aep\_core}'s workflow, and five for B0--B2, which
      have no window between writing a record and acknowledging it durable
      because they write no record.

\item \textbf{The prevention result no longer rests on one cell; it still rests
      on one host.} The original hard-Redis-kill ablation was collected on
      \textsc{no-readback} only, at one crash point, at \AepKillRuns{} runs per
      system, and was the whole of the barrier's measured case. The paper's
      prevention subsection now leads with a controlled-fault collection over
      \ArmASessions{} sessions and all \ArmAClasses{} capability classes, whose
      between-session spread is \ArmASpreadMax{} count per class. What has not
      changed is the machine: every collection here comes from one host, so the
      effect size's dependence on host timing is now measured \emph{on} that
      host rather than established \emph{across} hosts, and a second machine
      would still be sampling something we have not observed.

      We predicted that the other two capability classes would move the
      \emph{declared-ambiguity} column and not the applied-effect column,
      because whether an effect reached the provider is not something a
      read-back capability can change after the fact. We since collected those
      cells: \ClassSessions{} pre-registered sessions, \ClassRunsPerArm{} runs
      per arm per session, run-level interleaved, with the analysis fixed in
      advance. \textbf{The applied-effect column moved, and it moved in both
      directions.} The \textsc{authoritative-readback} minus
      \textsc{no-readback} applied-rate difference ran \ClassPpOne{},
      \ClassPpTwo{}, \ClassPpThree{} and \ClassPpFour{} percentage points across
      the \ClassSessions{} sessions --- a spread from \ClassPpMin{} to
      \ClassPpMax{} points on a comparison we expected to sit at zero. That is
      our own prediction contradicted in our own data, and we report it rather
      than resolving it.

      The session-clustered interval on that difference contains zero, at
      [\ClassPpLow{},\,\ClassPpHigh{}] percentage points. \textbf{It contains
      zero because the sessions disagree with one another, not because the
      effect is absent:} the half-width is \ClassPpHalfWidth{} points, wider
      than the mean it brackets at \ClassPpMean{} points, so this design could
      not have resolved the effect it observed. We registered a minimum
      detectable effect in advance and do not quote it here, because it was
      derived without a between-session variance component and so understates
      the precision this question needed; the comparison we make instead is the
      half-width against the observed mean. \textbf{What we report is therefore
      a failure to reject at a precision inadequate to the question --- not
      evidence that capability class leaves the applied-effect column alone.}
      The argument above is no longer unsupported: it is contradicted in the
      data, and untested at adequate power.

\item \textbf{The in-flight Redis-kill variant} (kill during transmission rather
      than before acknowledgement) \textbf{is now collected, and it ties.} Over
      \InflightSessions{} sessions the two arms differ by at most
      \InflightMaxDiff{} count in any capability class, both at ceiling: the tie
      our analysis predicted, measured rather than inferred. Two qualifications
      travel with it. The second session is a \emph{deterministic replay} of the
      first rather than an independent draw --- the harness assigns seeds per
      cell and repetition, so where no race remains the second collection
      reproduces the first exactly, which is itself the clearest evidence that
      none does. And B3's lowest cell is \InflightBthreeFloor{}/\AepKillRuns{},
      exactly the floor below which we would have read the result as a mis-timed
      injector rather than as a tie; it is the boundary, not headroom.

\item \textbf{The $30\%$ crash-probability regime} is implemented and not
      collected, so every rate in the paper comes from either an all-crash or a
      no-crash regime and nothing in between.

\item \textbf{The alternative read-back keying} is a sensitivity variant that
      has not been run; all results use one keying.

\item \textbf{One run did not complete} and is listed as such in the manifest
      rather than dropped silently.
\end{itemize}

A rate we did not measure is absent from the paper's tables rather than
interpolated, and the manifest shipped with the results states the run count for
every cell so a reader can see the shape of the coverage without reconstructing
it.

\section{Measurement incidents found during collection}
\label{supp:incidents}

The main paper's internal-validity subsection states each of these in one
sentence and points here. All three are defects in the \emph{measurement}, found
by the fault-injection matrix rather than by the unit suite, and all three are
fixed; they are recorded because the general point --- that a fault-injection
harness is a stronger verification tool than a unit suite --- cuts both ways,
and there may be defects it has not found.

\subsection{A resume defect that would have inflated counts}

Our matrix runner supports resumption, and a run whose previous attempt was
interrupted left event shards on disk that a fresh attempt could merge into its
own log, double-counting executions that were never part of the run being
recorded. It surfaced as a configuration-digest mismatch --- the digest check
did what it exists to do --- but between two attempts of the same harness
version the merge would have been \emph{silent}, and the error class is ``a
resumed run over-counts'', which for a paper measuring rates is as bad as errors
get. The safeguard is now explicit: the runner discards stale event shards and
any stale summary before a run writes anything, logs what it discarded, and
leaves the run's \emph{inputs} --- the rendered provider config and the oracle
ledger --- untouched, with a test pinning that distinction. Every resumed run in
the current results was collected after the fix; runs collected before it were
re-analysed and the single affected run was re-collected.

\subsection{Two further defects the matrix found that the unit suite had not}

A re-executing supervisor abandoned the lease instead of waiting for it, and
execution-state seeding was not idempotent. Both made a \emph{baseline} look
better than it is, so both were biasing against our own hypothesis, and both are
fixed.

\subsection{The test suite is destructive to a running matrix}

While collecting the B4/B4b cells we ran the project's own integration suite
against the same Redis. The suite is deliberately destructive: it deletes the
\texttt{aep:*} namespace in the database the matrix is using, and one of its
durability tests \texttt{docker restart}s the matrix's Redis container by name to
prove the AOF survives. Several runs died with connection errors, and the server
the batch was collecting against was replaced underneath it.

Runs that \emph{fail} are harmless --- they write no summary, and the resumption
logic re-runs them. The runs that \emph{completed} during that window are the
hazard, because a Redis blip makes a re-executing baseline retry, a baseline that
retries more produces more duplicates, and that bias points toward our own
hypothesis. They are also indistinguishable from clean runs in every artifact the
analysis reads. Every B4/B4b run collected in the disturbed window was therefore
discarded and re-collected on an otherwise idle host, and the server's
\texttt{run\_id} was recorded before and after the replacement batch to confirm
it was not restarted again; none of the discarded runs contributes a number to
the paper.

Two rules follow, and the artifact enforces both rather than documenting them:
the runner refuses to start beside a host-level fault injector, and it now
records the server's \texttt{run\_id} with every run so that a restart during
collection is visible in the results rather than inferred from a connection error
that happened to be logged. We report the episode because the generalisation is
not ``be careful'': it is that a fault-injection harness and a fault-injection
test suite make the same assumptions about owning the system, and running both is
a measurement error that produces ordinary-looking data.

\end{document}
